% Part: first-order-logic
% Chapter: natural-deduction
% Section: derivations

\documentclass[../../../include/open-logic-section]{subfiles}

\begin{document}

\iftag{FOL}
      {\olfileid{fol}{ntd}{der}}
      {\olfileid{pl}{ntd}{der}}

\olsection{اخذ (\usetoken{P}{derivation})}

\begin{explain}
اسیں دس چکے آں کہ مفروضہ کیہ ہوندا اے، تے استنباط دے قاعدے وی
دے چکے آں۔ فطری استنباط وچ !!^{derivation}s ایہناں توں
استقرائی طور تے بنائے جاندے نیں: ہر !!{derivation} یا تاں اکیلا
مفروضہ ہوندا اے، یا اک، دو، یا تِن !!{derivation}s دے بعد
درست استنباط لا کے بندا اے۔
\end{explain}

\begin{defn}[!!^{derivation}]
\Article{derivation} \emph{!!{derivation}}، !!a{sentence}~$!A$ دا
مفروضیاں~$\Gamma$ توں اک محدود درخت اے جیہڑا !!{sentence}s دا
بنیا ہووے تے ہِٹھ لکھیاں شرطاں پوری کردا ہووے:
\begin{enumerate}
\item درخت دے سب توں اُپرلے !!{sentence}s یا تاں $\Gamma$ وچ ہوندے
  نیں یا درخت دے کسے استنباط نال !!{discharged} ہوندے نیں۔
\item درخت دا سب توں ہِٹھلا !!{sentence}~$!A$ اے۔
\item درخت دا ہر !!{sentence}، ہِٹھلے آخرلے جملے~$!A$ توں سوا،
  استنباطی قاعدے دی درست ورتوں دا مقدمہ اے، تے اوس قاعدے دا نتیجہ
  درخت وچ اوس !!{sentence} دے ٹھیک ہِٹھاں آؤندا اے۔
\end{enumerate}
فیر اسیں $!A$ نوں !!{derivation} دا \emph{نتیجہ} تے $\Gamma$
نوں ایہدے !!{undischarged} مفروضے آکھدے آں۔

جے $!A$ دا~$\Gamma$ توں !!a{derivation} موجود ہووے، اسیں آکھدے
آں کہ $!A$،~$\Gamma$ توں \emph{!!{derivable}} اے، یا علامتاں وچ:
$\Gamma \Proves
!A$۔ جے~$!A$ دا !!a{derivation} ہووے جیہدے ہر مفروضے نوں
!!{discharged} کیتا گیا ہووے، اسیں~$\Proves !A$ لکھدے آں۔
\end{defn}

\begin{ex}
ہر مفروضہ اپنے آپ وچ !!a{derivation} اے۔ مثلاً $!A$ اپنے آپ وچ
!!a{derivation} اے، تے $!B$ وی۔ اسیں ایہناں توں، مثال لئی
$\Intro{\land}$ دا قاعدہ لا کے، نواں !!{derivation} بنا سکدے آں:
\begin{prooftree}
\AxiomC{$!A$}
\AxiomC{$!B$}
\RightLabel{\Intro{\land}}
\BinaryInfC{$!A \land !B$}
\end{prooftree}
ایہ قاعدے عام ورتوں لئی نیں: اسیں ایہدے وچ $!A$ تے~$!B$ دی
تھاں کوئی وی !!{sentence}s رکھ سکدے آں، مثلاً $!C$ تے~$!D$۔
فیر نتیجہ $!C \land !D$ ہووے گا، اس لئی
\begin{prooftree}
\AxiomC{$!C$}
\AxiomC{$!D$}
\RightLabel{\Intro{\land}}
\BinaryInfC{$!C \land !D$}
\end{prooftree}
اک درست !!{derivation} اے۔ یقیناً اسیں مفروضیاں دی ترتیب وی
بدل سکدے آں: $!D$،~$!A$ والا کردار تے $!C$،~$!B$ والا
کردار ادا کرے۔ ایس لئی
\begin{prooftree}
\AxiomC{$!D$}
\AxiomC{$!C$}
\RightLabel{\Intro{\land}}
\BinaryInfC{$!D \land !C$}
\end{prooftree}
وی اک درست !!{derivation} اے۔

ہن اسیں ہور قاعدہ، مثلاً $\Intro{\lif}$، لا سکدے آں۔ ایہ
شرطیہ نتیجہ کڈھن دی اجازت دیندا اے، تے اوہناں سارے مفروضیاں نوں
!!{discharge} کرن دی وی جیہڑے اوس شرطیہ دے مقدم دے عین برابر ہون۔
ایس لئی ہِٹھلے دونیں !!{derivation}s درست ہون گے:
\begin{prooftree}
\AxiomC{$\Discharge{!C}{1}$}
\AxiomC{$!D$}
\RightLabel{\Intro{\land}}
\BinaryInfC{$!C \land !D$}
\DischargeRule{\Intro{\lif}}{1}
\UnaryInfC{$!C \lif (!C \land !D)$}
\DisplayProof\bottomAlignProof
\AxiomC{$!C$}
\AxiomC{$\Discharge{!D}{1}$}
\RightLabel{\Intro{\land}}
\BinaryInfC{$!C \land !D$}
\DischargeRule{\Intro{\lif}}{1}
\UnaryInfC{$!D \lif (!C \land !D)$}
\end{prooftree}
ایہ بالترتیب وکھاندے نیں کہ $!D \Proves !C \lif (!C \land !D)$
تے $!C \Proves !D \lif (!C \land !D)$۔

یاد رکھو، مفروضے خارج کرنا اجازت اے، لازمی شرط نہیں: ساڈے لئی
مفروضے خارج کرنا ضروری نہیں۔ خاص طور تے، جے اوہ مفروضے
!!{derivation} وچ موجود وی نہ ہون، فیر وی اسیں قاعدہ لا سکدے آں۔
مثلاً ہِٹھلا استنباط جائز اے، بھاویں کسے مفروضے~$!A$ نوں
!!{discharged} کرن لئی موجود نہیں:
\begin{prooftree}
  \AxiomC{$!B$}
  \DischargeRule{\Intro{\lif}}{1}
  \UnaryInfC{$!A \lif !B$}
\end{prooftree}
\end{ex}

\end{document}
